\section*{Definitionen} 
\addcontentsline{toc}{section}{Definitionen}

\begin{acronym}[....]
	\acro{Hardstop}{Physicher Not-Aus über Pilztaster; verriegelt das System und setzt Motorcontroller sofort außer Betrieb}
	\acro{Softstop}{Kontrollierter Stopp des Fahrzeugs durch Taster oder \ac{UI}; keine Systemverriegelung}
	\acro{E-Stop}{Automatischer Notstopp durch Sensorik bei kritischem Hindernisabstand}
	\acro{Resume}{Wiederaufnahme des Fahrbetriebs durch automatische oder manuelle Freigabe}
	\acro{Kriechmodus}{Langsame, begrenzte Bewegung nach einem Nothalt zur sicheren Situationsklärung.}
	\acro{LED}{Light Emitting Diode}
	\acro{UI}{User Interface}
	\acro{PWM}{Pulse Width Modulation; Pulsweitenmodulation zur Steuerung von Motoren und Aktoren}
	\acro{ROS}{Robot Operating System; flexibles Framework für die Entwicklung von Robotersoftware}
	\acro{CAN}{Controller Area Network; serielles Bussystem zur Kommunikation zwischen Mikrocontrollern}
	\acro{LiDAR}{Light Detection and Ranging; optisches Messverfahren zur Entfernungsermittlung mittels Laserlicht}
	\acro{UML}{Unified Modeling Language}
	\acro{Warnmodus}{Blink-Signal mit Signalleuchten und keine Freigabe für Motorbewegung}
\end{acronym}


